% ============================================================
% Chapter 6 — Discussion (target 4–5 pp)
% ============================================================
\chapter{Discussion}
\label{ch:discussion}

\todo{Chapter not started. Skeleton below per the outline.}

\section{What Task Ordering Changes (and What It Does Not)}\label{sec:ordering}
% Interpretation of the 4 comparison axes.

\section{Parallel Task Coverage Revisited}\label{sec:ptc-revisited}
% What the evidence says about interactions, correlated samples,
% non-stationarity.

% --- §6.3–6.4 drafted 2026-07-14 (results-independent; re-read after
%     Ch. 5 is final to update the n-per-cell statement). ---
\section{Threats to Validity}
\label{sec:threats}

\paragraph{Simulator non-determinism under identical conditions.}
The dominant threat, quantified rather than assumed: three full-length
runs of byte-identical code, configuration and seed produced cumulative
vehicle success rates ranging 52.2--62.4\% (\cref{sec:iterations}).
CARLA under multi-agent load is not run-to-run deterministic even in
synchronous mode with all seeds fixed. Every effect in this thesis is
therefore read against this $\sim$10\,pp band: differences within it are
never attributed to an experimental factor, whatever their direction. The
band was measured on the pre-campaign trunk at training level; it is
assumed representative for campaign-era runs, which is itself an
approximation.

\paragraph{Sample size.}
Each cell contains a single run: the design was reduced from three paired
seeds per cell to one for compute reasons, a decision declared in
\cref{sec:campaign}. Every conclusion therefore rests on paired
contrasts, not on estimated distributions. Seed pairing and
byte-identical configurations make the contrasts internally clean, and
the primary contrast carries one built-in replication --- it is observed
independently under both critic architectures --- but no formal
statistical inference is possible at this $n$, and none is claimed. The
mitigations are the replicate band above and consistency checks ---
direction agreement across difficulty levels, evaluation scenarios,
agent classes, and measurement layers --- rather than significance
tests.

\paragraph{Evaluation noise.}
The final evaluation uses 100 episodes per scenario, i.e.\ 300 vehicle
records; at mid-range success rates the binomial standard error alone is
$\approx$3\,pp per scenario, before accounting for the within-episode
correlation of the three vehicles, which lowers the effective sample
further. Scenario-level differences of a few points are therefore not
individually meaningful; only patterns across scenarios are interpreted.

\paragraph{Reward--curriculum coupling.}
The reward function was tuned through gate iterations run mostly on
easy-locked configurations (\cref{sec:gate}), so the trunk's shaping may
be implicitly biased toward short routes. Both regimes share the
identical reward, so the primary contrast remains internally valid; the
caveat applies to interpreting \emph{absolute} performance on hard
routes, and to the possibility that a differently tuned reward would
shift the two regimes differently.

\paragraph{Single map pair.}
Generalization is measured on exactly one transfer, Town03 $\to$ Town05.
The two layouts differ in kind (\cref{sec:carla-bg}), which makes the
test meaningful, but claims are limited to this pair; no evidence is
offered that the ranking of regimes transfers to other maps, traffic
compositions, or weather conditions.

\paragraph{Measurement history.}
The evaluation pipeline itself failed validity once
(\cref{sec:bugs}), and the route-length defect silently redefined the
task distribution for a whole era of experiments. Both were caught by
integrity signals (implausibly regular outputs; missing per-agent logs).
The protocol now in place --- integrity checks, paired configurations,
recomputation from raw logs --- was built in response, but the general
lesson stands as a threat class: in a custom platform, apparent effects
must first survive the question ``is the instrument measuring what it
claims?''

\paragraph{Pedestrian route feasibility.}
Pedestrian results at medium/hard difficulty are confounded by a
structural environment limit: Town03's sidewalk topology often cannot
support the nominal route targets (\cref{sec:routes}). \Cref{sec:ped-limits}
separates environment-attributable from policy-attributable failure, and
pedestrian-side regime conclusions are restricted accordingly.

% ------------------------------------------------------------
\section{Scope Limitations}
\label{sec:scope}

Beyond threats to internal validity, the design deliberately fixes
several dimensions; findings should not be read beyond them.

\begin{itemize}
  \item \textbf{Actor architecture is fixed.} Both classes use MLP
        actors throughout; only the \emph{critic encoder} varies (MLP
        vs.\ GNN, training-side only). This thesis makes no claim about
        policy-network architecture, and the GNN axis speaks only to
        value estimation under CTDE.
  \item \textbf{The shared-world frame is not ablated.} Parallel task
        coverage (\cref{sec:parallel-coverage}) is the fixed frame of
        both arms, not a treatment: there is no 1-agent-per-world
        control, so its causal contribution to sample efficiency is
        characterized, not isolated.
  \item \textbf{Privileged observations, no perception.} Agents observe
        curated state features (\cref{sec:obs-actions}), not sensors.
        All findings concern decision-making under interaction;
        nothing follows for camera- or lidar-based pipelines.
  \item \textbf{One algorithm, one platform version.} MAPPO on
        RLlib~2.10.0 and CARLA~0.9.16 throughout. Regime effects are
        conditional on this stack; other algorithms may interact with
        task ordering differently.
  \item \textbf{Difficulty is route length only.} The configuration
        defines traffic-density and mixed difficulty families, but they
        are unused (\cref{sec:campaign}); ``difficulty'' in every claim
        means target route length at fixed traffic.
  \item \textbf{Fixed population.} Always 3 vehicles $+$ 3 pedestrians;
        no scaling evidence in either direction.
  \item \textbf{Pedestrian route generation.} The sidewalk-chain planner
        with its crossing cap is part of the platform being studied; its
        known feasibility limits bound what pedestrian policies can
        express, independently of training regime.
\end{itemize}
